% ID: SEL-D0021-OA-E02
% ESTADO: validada. S022, revisada con los criterios oficiales.
\textbf{a)} $B(100)=-100^2+360(100)-18000=\boxed{8000\text{ euros}}$.
Al completar el cuadrado, $B(x)=14400-(x-180)^2$.
Igualar a 13500 da $(x-180)^2=900$, así que
$\boxed{x=150\text{ o }x=210}$. Ambas cantidades pertenecen al dominio.

\textbf{b)} El cuadrado se anula sólo en $x=180$, por lo que
el máximo se alcanza al vender 180 unidades y vale $\boxed{14400\text{ euros}}$.

\textbf{c)} La gráfica es el arco de una parábola cóncava hacia abajo
para $50\le x\le350$. El vértice es $(180,14400)$ y las raíces
satisfacen $(x-180)^2=14400$, de donde $x=60$ y $x=300$.
Así, $B(x)\ge0$ para $\boxed{60\le x\le300}$, incluidos ambos umbrales.
\begin{center}
\setlength{\unitlength}{2mm}
\begin{picture}(43,34)(-5,-17)
\qbezier(-2,0)(16,0)(34,0)\put(35,0){$x$}
\qbezier(-2,-15)(-2,0)(-2,15)\put(-5,15){$B/1000$}
\put(0,1){\makebox(0,0)[r]{$50$}}
\qbezier(0,-2.5)(15,36.5)(30,-14.5)
\put(1,-2){\makebox(0,0)[t]{$60$}}
\put(25,1){\makebox(0,0)[b]{$300$}}
\put(13,15){\makebox(0,0)[b]{$(180,14.4)$}}
\put(30,-15){\makebox(0,0)[t]{$350$}}
\end{picture}
\end{center}
